\documentclass[12pt]{article}

% -------------------------------------------------
% Encoding, fonts, and typography
% -------------------------------------------------
\usepackage[T1]{fontenc}
\usepackage[utf8]{inputenc}
\usepackage{lmodern}
\usepackage{microtype}

% -------------------------------------------------
% Page layout
% -------------------------------------------------
\usepackage[
  margin=1in
]{geometry}

\usepackage{setspace}
\onehalfspacing

% -------------------------------------------------
% Mathematics
% -------------------------------------------------
\usepackage{amsmath}
\usepackage{amssymb}
\usepackage{amsthm}
\usepackage{mathtools}
\usepackage{bm}
\usepackage{bbm}

% -------------------------------------------------
% Tables
% -------------------------------------------------
\usepackage{booktabs}
\usepackage{threeparttable}
\usepackage{tabularx}
\usepackage{array}
\usepackage{multirow}
\usepackage{makecell}
\usepackage{siunitx}

\sisetup{
  detect-all,
  group-separator = {,},
  group-minimum-digits = 4
}

% -------------------------------------------------
% Figures
% -------------------------------------------------
\usepackage{graphicx}
\usepackage{subcaption}
\usepackage{float}

% -------------------------------------------------
% Lists and formatting
% -------------------------------------------------
\usepackage{enumitem}
\usepackage{xcolor}

% -------------------------------------------------
% References and citations
% -------------------------------------------------
\usepackage[
  authoryear,
  round
]{natbib}

\bibliographystyle{aer}

% -------------------------------------------------
% Hyperlinks and cross-references
% -------------------------------------------------
\usepackage[
  colorlinks=true,
  linkcolor=blue,
  citecolor=blue,
  urlcolor=blue
]{hyperref}

\usepackage[nameinlink,noabbrev]{cleveref}

% -------------------------------------------------
% Headers and footers
% -------------------------------------------------
\usepackage{fancyhdr}
\pagestyle{fancy}
\fancyhf{}
\fancyfoot[C]{\thepage}
\renewcommand{\headrulewidth}{0pt}

% -------------------------------------------------
% Section formatting
% -------------------------------------------------
\usepackage{titlesec}

\titleformat{\section}
  {\large\bfseries}
  {\thesection.}
  {0.5em}
  {}

\titleformat{\subsection}
  {\normalsize\bfseries}
  {\thesubsection.}
  {0.5em}
  {}

% -------------------------------------------------
% Theorem environments
% -------------------------------------------------
\newtheorem{proposition}{Proposition}
\newtheorem{lemma}{Lemma}
\newtheorem{corollary}{Corollary}
\newtheorem{assumption}{Assumption}

\theoremstyle{definition}
\newtheorem{definition}{Definition}

\theoremstyle{remark}
\newtheorem{remark}{Remark}

% -------------------------------------------------
% Common economics/math operators
% -------------------------------------------------
\DeclareMathOperator{\E}{\mathbb{E}}
\DeclareMathOperator{\Var}{Var}
\DeclareMathOperator{\Cov}{Cov}
\DeclareMathOperator{\Corr}{Corr}
\DeclareMathOperator{\argmax}{arg\,max}
\DeclareMathOperator{\argmin}{arg\,min}

% -------------------------------------------------
% Common shortcuts
% -------------------------------------------------
\newcommand{\R}{\mathbb{R}}
\newcommand{\ind}{\mathbbm{1}}
\newcommand{\eps}{\varepsilon}
\newcommand{\dd}{\,\mathrm{d}}

% -------------------------------------------------
% Equation numbering
% -------------------------------------------------
\numberwithin{equation}{section}

% -------------------------------------------------
% Paragraph formatting
% -------------------------------------------------
\setlength{\parindent}{1.5em}
\setlength{\parskip}{0pt}

% -------------------------------------------------
% Document
% -------------------------------------------------
\begin{document}

\section{Known type masses and timing}
There are $N$ consumers, two inside goods $j\in\{1,2\}$ owned by a
monopolist, and an outside option $0$. Types are $r\in\{H,L\}$, with known,
fixed probability masses
\[
 \omega_H=\lambda,\qquad \omega_L=1-\lambda,\qquad
 0\leq\lambda\leq1,\qquad \omega_H+\omega_L=1.
\]
Individual type is unobserved in the main common-price problem. The firm
chooses its entire price menu before purchases occur, using these known masses.
There is no learning, estimation, posterior, or subsequent belief-based
repricing. A terminal menu may depend on the purchased product because that
product determines the next period's variety utility; this is a preannounced
pricing rule. Commitment makes this timing explicit.

Consumers are myopic, while the firm discounts second-period profits by
$\delta$. Consumers begin with no prior purchase. First-period utility has
no variety bonus; variety utility begins in period two.
An outside purchase in period one exits permanently. In period two, $h$
denotes the inside product purchased in period one. All prices are net
transaction prices in $[\underline p,\overline p]=[1,8]$.

\section{Demand: every differentiation step}
At a history $h$, define
\begin{align}
 V_{rj}(h;p)&=\beta_j+\gamma_r\ind\{j\ne h\}+\alpha p_j,
 &\alpha&<0,\\
 Z_r(h;p)&=1+\sum_{k=1}^2 e^{V_{rk}(h;p)},\\
 s_{rj}(h;p)&=\frac{e^{V_{rj}(h;p)}}{Z_r(h;p)},
 &s_{r0}(h;p)&=\frac1{Z_r(h;p)},\\
 s_{r0}+s_{r1}+s_{r2}
 &=\frac{1+e^{V_{r1}}+e^{V_{r2}}}{Z_r}=1.
\end{align}
Suppress history and prices temporarily. For a price $p_j$,
\begin{align}
 \frac{\partial V_{rk}}{\partial p_j}&=\alpha\ind\{k=j\},\\
 \frac{\partial e^{V_{rk}}}{\partial p_j}
 &=e^{V_{rk}}\frac{\partial V_{rk}}{\partial p_j}
 =\alpha e^{V_{rk}}\ind\{k=j\},\\
 \frac{\partial Z_r}{\partial p_j}
 &=\sum_{l=1}^2\alpha e^{V_{rl}}\ind\{l=j\}
 =\alpha e^{V_{rj}},\\
 \frac{\partial s_{rk}}{\partial p_j}
 &=\frac{Z_r\alpha e^{V_{rk}}\ind\{k=j\}
               -e^{V_{rk}}\alpha e^{V_{rj}}}{Z_r^2}\\
 &=\alpha\frac{e^{V_{rk}}}{Z_r}\ind\{k=j\}
       -\alpha\frac{e^{V_{rk}}}{Z_r}\frac{e^{V_{rj}}}{Z_r}\\
 &=\alpha s_{rk}(\ind\{k=j\}-s_{rj}).
\end{align}
Thus the own and cross elasticities are, respectively,
\begin{align}
 \epsilon_{rjj}
 &=\frac{p_j}{s_{rj}}\frac{\partial s_{rj}}{\partial p_j}
 =\frac{p_j}{s_{rj}}\alpha s_{rj}(1-s_{rj})
 =\alpha p_j(1-s_{rj}),\\
 \epsilon_{rkj}
 &=\frac{p_j}{s_{rk}}\frac{\partial s_{rk}}{\partial p_j}
 =\frac{p_j}{s_{rk}}(-\alpha s_{rk}s_{rj})
 =-\alpha p_js_{rj},\qquad k\ne j.
\end{align}
First-period utility and shares have no purchase-history state:
\begin{align}
 V_{rj1}&=\beta_j+\alpha p_{j1},\\
 S_{rj}(p_1)&=\frac{e^{\beta_j+\alpha p_{j1}}}
 {1+\sum_k e^{\beta_k+\alpha p_{k1}}},\\
 S_{Hj}&=S_{Lj},\\
 D_{r,kj}(p_1):=\frac{\partial S_{rk}}{\partial p_{j1}}
 &=\alpha S_{rk}(\ind\{k=j\}-S_{rj}).
\end{align}
The first purchase has the same demand for both types. Thus purchase masses
$\omega_r S_{rh}$ factor into the known type mass and a common purchase rate.

\section{Common prices chosen using the known masses}
Write $p_{k2}(h)$ for the terminal price of product $k$ following product $h$.
For any proposed menu, define
\begin{align}
 m_{j1}&=p_{j1}-c_j,& m_{k2}(h)&=p_{k2}(h)-c_k,\\
 w_r(h;p_2(h))&=\sum_{k=1}^2s_{rk}(h;p_2(h))m_{k2}(h),\\
 A_{rh}&=m_{h1}+\delta w_r(h;p_2(h)).
\end{align}
The expected mass of type-$r$ consumers purchasing $h$ is
$N\omega_rS_{rh}$. This is a joint purchase mass, computed before purchases,
not an updated type probability. Expected profits follow by multiplying each
mass by its associated margin:
\begin{align}
 \pi_1&=N\sum_r\sum_h\omega_rS_{rh}m_{h1},\\
 \pi_2&=N\sum_r\sum_h\sum_k\omega_rS_{rh}
                   s_{rk}(h;p_2(h))m_{k2}(h),\\
 \Pi&=\pi_1+\delta\pi_2\\
 &=N\sum_r\sum_h\omega_rS_{rh}m_{h1}
    +N\delta\sum_r\sum_h\omega_rS_{rh}
                         \sum_k s_{rk}(h;p_2(h))m_{k2}(h)\\
 &=N\sum_r\omega_r\sum_hS_{rh}
                  [m_{h1}+\delta w_r(h;p_2(h))]\\
 &=N\sum_r\omega_r\sum_hS_{rh}A_{rh}.
\end{align}
The common-menu problem is
\[
 \max_{p_1,p_2(1),p_2(2)\in[1,8]^2}\Pi.
\]
All six prices are selected together. ``Common'' means across types; products,
periods, and purchase histories may have different prices.

\subsection{First-period derivatives}
For a partial derivative, hold the other menu entries fixed. In particular,
$w_r(h;p_2(h))$ has no direct dependence on $p_{j1}$. Therefore
\begin{align}
 \frac{\partial m_{h1}}{\partial p_{j1}}&=\ind\{h=j\},\\
 \frac{\partial A_{rh}}{\partial p_{j1}}
 &=\ind\{h=j\}+\delta\,0=\ind\{h=j\},\\
 \frac1N\frac{\partial\Pi}{\partial p_{j1}}
 &=\sum_r\omega_r\sum_h
       \frac{\partial(S_{rh}A_{rh})}{\partial p_{j1}}\\
 &=\sum_r\omega_r\sum_h
       \left[\frac{\partial S_{rh}}{\partial p_{j1}}A_{rh}
                   +S_{rh}\frac{\partial A_{rh}}{\partial p_{j1}}\right]\\
 &=\sum_r\omega_r\left[\sum_hD_{r,hj}A_{rh}
                          +\sum_hS_{rh}\ind\{h=j\}\right]\\
 &=\sum_r\omega_r\left[S_{rj}+\sum_hD_{r,hj}
                         (m_{h1}+\delta w_r(h;p_2(h)))\right].
\end{align}
At an interior optimum this equals zero. With two types it is explicitly
\begin{align}
 0={}&\lambda\left[S_{Hj}+D_{H,1j}A_{H1}+D_{H,2j}A_{H2}\right]\notag\\
 &+(1-\lambda)\left[S_{Lj}+D_{L,1j}A_{L1}+D_{L,2j}A_{L2}\right].
\end{align}
No derivative of a type probability appears: $\lambda$ is fixed.

\subsection{Second-period derivatives}
For a given history $h$, first differentiate conditional terminal profit:
\begin{align}
 \frac{\partial w_r(h)}{\partial p_{j2}(h)}
 &=\sum_k\left[\frac{\partial s_{rk}(h)}{\partial p_{j2}(h)}m_{k2}(h)
                      +s_{rk}(h)\ind\{k=j\}\right]\\
 &=\sum_k\alpha s_{rk}(h)(\ind\{k=j\}-s_{rj}(h))m_{k2}(h)
          +s_{rj}(h)\\
 &=\alpha s_{rj}(h)m_{j2}(h)
       -\alpha s_{rj}(h)\sum_ks_{rk}(h)m_{k2}(h)+s_{rj}(h)\\
 &=s_{rj}(h)\left[1+\alpha(m_{j2}(h)-w_r(h))\right].
\end{align}
Only continuation at history $h$ uses the entry $p_{j2}(h)$, so
\begin{align}
 \frac1N\frac{\partial\Pi}{\partial p_{j2}(h)}
 &=\delta\sum_r\omega_rS_{rh}\frac{\partial w_r(h)}{\partial p_{j2}(h)}\\
 &=\delta\sum_r\omega_rS_{rh}s_{rj}(h)
                   [1+\alpha(m_{j2}(h)-w_r(h))]\\
 &=\delta\{\lambda S_{Hh}s_{Hj}(h)[1+\alpha(m_{j2}(h)-w_H(h))]\notag\\
 &\hspace{2.4em}+(1-\lambda)S_{Lh}s_{Lj}(h)
                         [1+\alpha(m_{j2}(h)-w_L(h))]\}.
\end{align}
Since $S_{Hh}=S_{Lh}$, the common first-period purchase rate factors out
of each terminal condition. Terminal prices therefore maximize
$\sum_r\omega_r w_r(h)$ directly using the known type masses.
Equal terminal markups across goods need not hold for a mixture of types.

For any price coordinate $x$, write $g_x=\partial\Pi/\partial x$.
The necessary bound conditions for maximization are
\[
 g_x=0\ \text{if }1<x<8,\qquad
 g_x\leq0\ \text{if }x=1,\qquad
 g_x\geq0\ \text{if }x=8.
\]
At the lower bound only increases are feasible, so a positive derivative
would improve profit; at the upper bound only decreases are feasible, so a
negative derivative would improve profit. These conditions are necessary,
not a proof of a global maximum.

\section{Observed-type benchmark and price identities}
If type is observed exogenously, give each type a separate menu. Then
\begin{align}
 \Pi^{\mathrm{obs}}&=N\sum_r\omega_r\Pi_r,\\
 \Pi_r&=\sum_hS_{rh}[p_{rh1}-c_h+\delta w_r(h)],\\
 w_r(h)&=\max_{p_{r2}(h)\in[1,8]^2}
               \sum_ks_{rk}(h;p_{r2}(h))[p_{rk2}(h)-c_k].
\end{align}
For $\omega_r>0$, multiplying $\Pi_r$ by $N\omega_r$ does not change its
maximizing menu, because it multiplies every profit comparison by a positive
constant. At zero mass its menu is irrelevant to aggregate profit; the code
still computes it as a benchmark. These observed-type menus have no learning.
Terminal prices can be solved first, followed by the initial prices.

The initial derivative is the common-price derivative with a single type:
\[
 0=S_{rj}+\sum_hD_{r,hj}[m_{rh1}+\delta w_r(h)].
\]
Write $s_{rk1}=S_{rk}$ and
$a_{rk}=m_{rk1}+\delta w_r(k)$. Substitution and division by $s_{rj1}>0$ give
\begin{align}
 0&=s_{rj1}+\sum_k\alpha s_{rk1}(\ind\{k=j\}-s_{rj1})a_{rk}\\
 &=s_{rj1}+\alpha s_{rj1}a_{rj}
                         -\alpha s_{rj1}\sum_ks_{rk1}a_{rk},\\
 0&=1+\alpha\left[a_{rj}-\sum_ks_{rk1}a_{rk}\right].
\end{align}
To display the outside option explicitly,
\begin{align}
 a_{rj}-\sum_ks_{rk1}a_{rk}
 &=(1-s_{rj1})a_{rj}-\sum_{k\ne j}s_{rk1}a_{rk}\\
 &=\left(s_{r01}+\sum_{k\ne j}s_{rk1}\right)a_{rj}
                              -\sum_{k\ne j}s_{rk1}a_{rk}\\
 &=s_{r01}a_{rj}+\sum_{k\ne j}s_{rk1}(a_{rj}-a_{rk}).
\end{align}
Solving the first-order equation algebraically,
\begin{align}
 -1&=\alpha\left[a_{rj}-\sum_ks_{rk1}a_{rk}\right],\\
 -\frac1\alpha&=a_{rj}-\sum_ks_{rk1}a_{rk},\\
 a_{rj}&=-\frac1\alpha+\sum_ks_{rk1}a_{rk},\\
 p_{rj1}-c_j+\delta w_r(j)
 &=-\frac1\alpha+\sum_ks_{rk1}[p_{rk1}-c_k+\delta w_r(k)],\\
 p_{rj1}&=c_j-\delta w_r(j)-\frac1\alpha
                  +\sum_ks_{rk1}[p_{rk1}-c_k+\delta w_r(k)].
\end{align}
Subtracting the equations for $j$ and $l$ cancels the common sum:
\begin{align}
 a_{rj}-a_{rl}&=0,\\
 (p_{rj1}-c_j+\delta w_r(j))-(p_{rl1}-c_l+\delta w_r(l))&=0,\\
 p_{rj1}-p_{rl1}&=c_j-c_l-\delta[w_r(j)-w_r(l)].
\end{align}
These identities are implicit, since shares depend on prices, and this
relative-price identity assumes interior prices.

For a single observed type at terminal history $h$, the terminal derivative
already derived gives
\begin{align}
 0&=s_{rj2}(h)[1+\alpha(m_{rj2}(h)-w_r(h))],\\
 0&=1+\alpha(m_{rj2}(h)-w_r(h)),\\
 m_{rj2}(h)&=w_r(h)-1/\alpha.
\end{align}
The right side is the same for every $j$, so all inside margins equal $M_r(h)$.
Consequently
\begin{align}
 w_r(h)&=\sum_ks_{rk2}(h)M_r(h)
       =(1-s_{r02}(h))M_r(h),\\
 M_r(h)&=(1-s_{r02}(h))M_r(h)-1/\alpha,\\
 M_r(h)-(1-s_{r02}(h))M_r(h)&=-1/\alpha,\\
 s_{r02}(h)M_r(h)&=-1/\alpha,\\
 M_r(h)&=-\frac1{\alpha s_{r02}(h)},\\
 p_{rj2}(h)-p_{rl2}(h)
 &=[c_j+M_r(h)]-[c_l+M_r(h)]=c_j-c_l.
\end{align}
The outside share must be retained. These equalities assume interior prices.

\section{Pricing directly from a known two-point distribution}
All consumers begin with no previous product. Write $\rho=\Pr(L)=1-\lambda$, so
$\Pr(H)=1-\rho$. There is no first-period switching bonus:
\begin{align}
 V_{rj1}^{B}&=\beta_j+\alpha p_{j1},\\
 b_j(p_1)&=\frac{\exp(\beta_j+\alpha p_{j1})}
                    {1+\sum_k\exp(\beta_k+\alpha p_{k1})},\\
 S_{Hj}^{B}&=S_{Lj}^{B}=b_j(p_1).
\end{align}
Known type masses are used directly in terminal expected profits:
\begin{align}
 \bar w(h;p)&=(1-\rho)w_H(h;p)+\rho w_L(h;p),\\
 p_2^*(h;\rho)&\in\argmax_{p\in[1,8]^2}\bar w(h;p),\\
 W_h(\rho)&=\bar w(h;p_2^*(h;\rho)).
\end{align}
The terminal interior condition is explicitly
\begin{align}
 0=\frac{\partial\bar w(h;p)}{\partial p_j}
 ={}&(1-\rho)s_{Hj}(h;p)[1+\alpha(p_j-c_j-w_H(h;p))]\notag\\
 &+\rho s_{Lj}(h;p)[1+\alpha(p_j-c_j-w_L(h;p))].
\end{align}
Since first-period shares are identical across types, the total objective
factors without any conditioning on type:
\begin{align}
 \Pi^{B}/N
 &=(1-\rho)\sum_hb_h[m_{h1}+\delta w_H(h)]
       +\rho\sum_hb_h[m_{h1}+\delta w_L(h)]\\
 &=\sum_hb_h[(1-\rho)m_{h1}+\rho m_{h1}
                 +\delta((1-\rho)w_H(h)+\rho w_L(h))]\\
 &=\sum_hb_h[m_{h1}+\delta\bar w(h)]\\
 &=\sum_hb_h[p_{h1}-c_h+\delta W_h(\rho)]
       \quad\text{at the optimal terminal menu}.
\end{align}
Because $W_h(\rho)$ has no dependence on $p_1$, terminal prices can be chosen
first. Put $A_h^B=p_{h1}-c_h+\delta W_h(\rho)$. Then
\begin{align}
 \frac1N\frac{\partial\Pi^B}{\partial p_{j1}}
 &=\sum_h\left[\alpha b_h(\ind\{h=j\}-b_j)A_h^B
                                  +b_h\ind\{h=j\}\right]\\
 &=\alpha b_jA_j^B-\alpha b_j\sum_hb_hA_h^B+b_j\\
 &=b_j\left[1+\alpha\left(A_j^B-\sum_hb_hA_h^B\right)\right].
\end{align}
For an interior solution, divide by $b_j>0$ and rearrange:
\begin{align}
 0&=1+\alpha\left(A_j^B-\sum_hb_hA_h^B\right),\\
 A_j^B&=-1/\alpha+\sum_hb_hA_h^B,\\
 p_{j1}&=c_j-\delta W_j(\rho)-1/\alpha
                   +\sum_hb_h[p_{h1}-c_h+\delta W_h(\rho)].
\end{align}
This is pricing under a known probability mass function, with no learning.

\subsection{A fixed discount decision and every cutoff step}
Construct a regular menu $R$ from the known-H newborn optimum. Construct menu
$D$ by holding its initial prices fixed and replacing terminal prices with the
known-L terminal optimum. The code verifies that these terminal prices are
weakly lower. The full known-L dynamic policy is not necessarily a discount:
its introductory prices can be higher because the H consumer has greater
continuation value. Both comparison menus stay fixed as $\rho$ varies.

For $X\in\{R,D\}$, define profits conditional on true type and the profit
change from discounting:
\begin{align}
 F_r^X&=\sum_h b_h(p_1^R)
      [p_{h1}^R-c_h+\delta w_r(h;p_2^X(h))],\\
 \Delta_r&=F_r^D-F_r^R\\
 &=\sum_hb_h(p_1^R)\{[p_{h1}^R-c_h+\delta w_r(h;p_2^D(h))]
                         -[p_{h1}^R-c_h+\delta w_r(h;p_2^R(h))]\}\\
 &=\delta\sum_hb_h(p_1^R)
                         [w_r(h;p_2^D(h))-w_r(h;p_2^R(h))].
\end{align}
By the known type probabilities, the expected gain is
\begin{align}
 B(\rho)&=[(1-\rho)F_H^D+\rho F_L^D]
                         -[(1-\rho)F_H^R+\rho F_L^R]\\
 &=(1-\rho)(F_H^D-F_H^R)+\rho(F_L^D-F_L^R)\\
 &=(1-\rho)\Delta_H+\rho\Delta_L\\
 &=\Delta_H-\rho\Delta_H+\rho\Delta_L\\
 &=\Delta_H+\rho(\Delta_L-\Delta_H).
\end{align}
For $\Delta_H<0<\Delta_L$, the slope is positive, and the indifference point
satisfies
\begin{align}
 0&=\Delta_H+\rho^*(\Delta_L-\Delta_H),\\
 -\Delta_H&=\rho^*(\Delta_L-\Delta_H),\\
 \rho^*&=\frac{-\Delta_H}{\Delta_L-\Delta_H}.
\end{align}
Because the denominator is positive, the discount rule is obtained without
reversing the inequality:
\begin{align}
 B(\rho)>0
 &\iff\Delta_H+\rho(\Delta_L-\Delta_H)>0\\
 &\iff\rho(\Delta_L-\Delta_H)>-\Delta_H\\
 &\iff\rho>\frac{-\Delta_H}{\Delta_L-\Delta_H}=\rho^*.
\end{align}
For these menus the numerical values are approximately
\begin{align}
 \Delta_H&=-0.04484198,&\Delta_L&=0.02410654,\\
 \Delta_L-\Delta_H&=0.02410654-(-0.04484198)=0.06894852,\\
 \rho^*&\simeq\frac{0.04484198}{0.06894852}=0.650369,\\
 B(1/2)&\simeq\tfrac12(-0.04484198)+\tfrac12(0.02410654)\\
 &=-0.02242099+0.01205327=-0.01036772,\\
 1-\rho^*&\simeq1-0.650369=0.349631.
\end{align}
Offer the fixed discount above approximately $65.04\%$ probability of L;
below it retain the regular menu; at equality either menu is optimal.
At a 50/50 mass the discount loses profit. The cost of discounting an H type
is included in $\Delta_H$ and requires no extra fee.

If $a\in[0,1]$ instead denotes the probability of sending the fixed discount,
\begin{align}
 F(a;\rho)&=(1-a)F^R(\rho)+aF^D(\rho)\\
 &=F^R(\rho)-aF^R(\rho)+aF^D(\rho)\\
 &=F^R(\rho)+aB(\rho),\\
 \frac{\partial F(a;\rho)}{\partial a}&=B(\rho).
\end{align}
Thus $a=1$ maximizes profit when $B>0$, $a=0$ when $B<0$, and every $a$
is optimal when $B=0$. This is distinct from varying the type mass $\rho$.
A freely optimized common menu weakly dominates either fixed menu because
both fixed menus are feasible choices in its optimization problem.

\section{Numerical implementation and profit distribution}
The parameters are
\[
 \beta=(2.2,2.0),\quad c=(1.2,1.0),\quad\alpha=-1.5,\quad
 \delta=0.95,\quad\gamma_H=2,\quad\gamma_L=0.1.
\]
Run \texttt{python3 toy\_model.py}. The common-menu solver jointly
chooses all prices using the fixed mass objective above. Terminal problems
can also be solved first because first-period demand is identical across types.
Optimization uses three starting values and checked solver status; this is not
a global optimality proof. The fixed terminal-discount comparison uses the same
model and holds the regular first-period prices fixed.

The results directory contains \path{prices_50_50.csv},
\path{population_mix.csv}, \path{profit_by_gamma.csv},
\path{newborn_prices.csv}, and
\path{newborn_discount_probability.csv}, together
with reports and figures. The newborn CSV labels $\rho$ as
\texttt{probability\_L}. The population experiment varies $\lambda=\Pr(H)$;
the newborn experiment varies $\rho=\Pr(L)$.

For any fixed common policy $P$, let $f_r(P)$ be expected discounted profit
conditional on type. The distribution of these conditional expectations has
two masses, with cumulative distribution function
\[
 \Pr(f_r(P)\leq x)
 =\lambda\ind\{f_H(P)\leq x\}
       +(1-\lambda)\ind\{f_L(P)\leq x\}.
\]
Its expectation is
\[
 \E[f_r(P)]=\lambda f_H(P)+(1-\lambda)f_L(P).
\]
If both profit values coincide, their masses add to one at that value.
This describes type-conditional expected profits, rather than realized
transaction profits. A curve of expected profit against $\gamma$ is not a
probability density; a continuous type distribution would have to be specified
to construct a continuous distribution of conditional expected profits.

\end{document}
